\section{Formative review of artifacts}\label{FormativeStudy}

Head et al.'s prior content analysis~\cite{Head2022-ws} identified four patterns of interactive formula augmentation: scrubbable expressions, editable text fields, linked selections, and external controls. Their focus was on augmentations broadly. Given our focus on interaction, we sought additional documents with this pattern to identify common augmentation forms and better understand potential frictions in their implementation.

We acquired 11 documents containing formulas of interest from sources including Distill.pub, Setosa.io, Explorable Explanations, complex-analysis.com, Georgia Tech's Interactive Linear Algebra textbook, and personal websites. Sources were curated from \citet{Head2022-ws}'s content analysis, \url{https://explorabl.es/}, \url{https://github.com/ubavic/awesome-interactive-math}, and our own web searches. Two authors analyzed 46 formulas of interest. All included formulas were either themselves dynamic or adjacent to and significant in an interactive mathematical explanation.

\paragraph{Augmentation patterns}
Of these, 24 were dynamic (changing following interaction) and 22 supported linking between formula and visualization elements. Formulas were controlled by sliders (8), draggable graph points (17), radio buttons (8), in-formula scrubbing (5), in-formula editing (1). Walkthroughs---step-by-step navigation through staged derivations---appeared for formulas in 1 example and for broader mathematical content in 2 others. We further reviewed the code from open source repositories we found for 11 documents and identified recurring implementation patterns representing potential frictions (see also appendix Table~\ref{tab:formative-audit}):

\begin{itemize}
\item \emph{F1: Diffuseness and tangling (11 articles).} A single interaction is implemented across multiple files; or rendering, layout, and behavior is tangled within a single function.
\item \emph{F2: Indirect identifiers (11 articles).} Variables are referred to in the code using identifiers or array indexing distinct from how the variable appears in the formula.
\item \emph{F3: Fussing the DOM and LaTeX (9 articles)}: The code constructed \LaTeX{} strings programmatically, traversed rendered formula DOM, or both. The two exceptions bypassed LaTeX with their own rendering engines.
\item\emph{F4: Repetition (10 articles)}: interaction logic appeared duplicated across instances of related math objects.
\item\emph{F5: State management (10 articles)}: Articles implemented custom state management and wired up drag handlers, click listeners, and other callbacks to propagate changes across formulas, variables, and other representations.
\end{itemize}

These frictions may distance authors from focusing on presenting a formula well. Our controlled study (Section~\ref{results}) compares Delta to a baseline assess whether improved tooling can reduce this work.